% !TEX TS-program = xelatex
\documentclass[10pt,a4paper]{article}

% Short CV | Editorial layout, aligned section rows
% Compile with XeLaTeX (Overleaf: Menu > Compiler > XeLaTeX)
\usepackage[a4paper,top=13mm,bottom=12mm,left=18mm,right=18mm]{geometry}
\usepackage{fontspec}
\usepackage{microtype}
\usepackage{xcolor}
\usepackage{hyperref}
\usepackage{enumitem}
\usepackage{ragged2e}
\usepackage{tikz}
\usepackage{array}
\usepackage{calc}
\usepackage{fancyhdr}

\setmainfont{LibertinusSerif-Regular.otf}[BoldFont=LibertinusSerif-Bold.otf,ItalicFont=LibertinusSerif-Italic.otf,BoldItalicFont=LibertinusSerif-BoldItalic.otf]
\setsansfont{Lato}
\definecolor{ink}{HTML}{172B3B}
\definecolor{deepblue}{HTML}{193C58}
\definecolor{blue}{HTML}{446982}
\definecolor{quiet}{HTML}{596C7B}
\definecolor{pale}{HTML}{D4DDE3}
\definecolor{warm}{HTML}{F2F5F6}
\definecolor{teal}{HTML}{4D7F86}
\definecolor{ochre}{HTML}{AC7855}
\definecolor{softteal}{HTML}{E7EFF0}

\hypersetup{
  colorlinks=true,
  urlcolor=teal,
  linkcolor=deepblue,
  pdfauthor={Adrián Vidal Delgado},
  pdftitle={Adrián Vidal Delgado | Short CV}
}
\urlstyle{same}
\pagestyle{empty}
\setlength{\parindent}{0pt}
\setlength{\parskip}{0pt}
\color{ink}
\linespread{1.035}
\hyphenpenalty=10000
\exhyphenpenalty=10000
\setlength{\emergencystretch}{2em}

% Type styles
\newcommand{\caps}[1]{{\sffamily\fontsize{8}{10}\selectfont\bfseries\addfontfeatures{LetterSpace=1.35}#1}}
\newcommand{\tinycaps}[1]{{\sffamily\fontsize{8}{10}\selectfont\color{quiet}\addfontfeatures{LetterSpace=0.9}#1}}
\newcommand{\subtle}[1]{{\sffamily\fontsize{9.1}{12.5}\selectfont\color{quiet}#1}}
\newcommand{\bodytext}[1]{{\sffamily\fontsize{9.25}{12.7}\selectfont#1}}
\newcommand{\sect}[1]{%
  \vspace{9pt}%
  {\caps{\textcolor{deepblue}{#1}}}\par\vspace{3.5pt}%
  {\color{pale}\rule{\linewidth}{0.5pt}}\par\vspace{6pt}%
}
\newcommand{\sectleft}[1]{%
  \vspace{18pt}%
  {\caps{\textcolor{deepblue}{#1}}}\par\vspace{4.5pt}%
  {\color{pale}\rule{\linewidth}{0.5pt}}\par\vspace{10pt}%
}
\newcommand{\asideentry}[3]{%
  {\sffamily\fontsize{9.2}{12.6}\selectfont\bfseries #1}\par
  {\sffamily\fontsize{8.55}{11.6}\selectfont\color{quiet}#2}\par
  \if\relax\detokenize{#3}\relax\else
    {\sffamily\fontsize{8.55}{11.5}\selectfont #3}\par
  \fi
  \vspace{8.4pt}%
}
% Each entry in the wide column is anchored by a slim, colored side tab.
\newcommand{\entrytab}[3]{%
  \noindent
  \begin{minipage}[t]{5pt}%
    \vspace{3.5pt}\color{#1}\rule{2.5pt}{14pt}%
  \end{minipage}\hspace{7.5pt}%
  \begin{minipage}[t]{\dimexpr\linewidth-12.5pt\relax}%
    \vspace{0pt}#2%
  \end{minipage}\par\vspace{#3}%
}
\newcommand{\paper}[3]{%
  \entrytab{ochre}{%
    {\fontsize{11.2}{14.2}\selectfont\bfseries\color{ink}#1}\par\vspace{2.3pt}%
    {\bodytext{#2}}\par\vspace{2pt}%
    {\sffamily\fontsize{8.9}{12}\selectfont\color{quiet}#3}\par
  }{11pt}%
}
\newcommand{\talk}[3]{%
  \entrytab{teal}{%
    {\sffamily\fontsize{8}{10}\selectfont\bfseries\addfontfeatures{LetterSpace=0.8}\color{ochre}#1}\par\vspace{2.1pt}%
    {\sffamily\fontsize{9.15}{12.5}\selectfont\bfseries #2}\par\vspace{1.5pt}%
    {\sffamily\fontsize{8.65}{11.7}\selectfont\color{quiet}#3}\par
  }{7pt}%
}
\newcommand{\lineentry}[2]{%
  \entrytab{blue}{%
    {\sffamily\fontsize{9.15}{12.35}\selectfont\textbf{#1}}\par
    {\sffamily\fontsize{8.75}{11.9}\selectfont\color{quiet}#2}\par
  }{6pt}%
}



% Paired rows align section headers across both columns and automatically
% continue at the height of the taller section. Edit each pair together.
\renewcommand{\sectleft}[1]{\sect{#1}}
\newcommand{\pairedsections}[2]{%
  \noindent\begin{tabular}{@{}p{0.325\textwidth}@{\hspace{0.032\textwidth}{\color{pale}\vrule width 0.65pt}\hspace{0.032\textwidth}}p{0.607\textwidth}@{}}
    \RaggedRight #1 & \RaggedRight #2
  \end{tabular}\par
}

\begin{document}

% ----------- Identity -----------
\begin{minipage}[t]{0.69\textwidth}
  {\sffamily\fontsize{8}{11}\selectfont\color{ochre}\bfseries\addfontfeatures{LetterSpace=1.8}SHORT CV}\par\vspace{8pt}
  {\fontsize{28}{30}\selectfont\color{deepblue}Adrián\par Vidal Delgado}\par\vspace{7pt}
  {\sffamily\fontsize{10.8}{13}\selectfont\color{ink}\bfseries PhD Candidate in Mathematics}\par
  {\sffamily\fontsize{9.3}{13}\selectfont\color{quiet}Institute of Mathematics · Universidad de Talca, Chile}
\end{minipage}%
\hfill
\begin{minipage}[t]{0.29\textwidth}
  \raggedleft
  \vspace{11pt}
  {\sffamily\fontsize{9.05}{13}\selectfont
  \href{mailto:adrian.vidal@utalca.cl}{adrian.vidal@utalca.cl}\par\vspace{4pt}
  \href{https://adrian279vd.github.io/}{adrian279vd.github.io}\par\vspace{4pt}
  \href{https://scholar.google.com/citations?user=HVv9miEAAAAJ\&hl=en}{Google Scholar}\; /\; \href{https://orcid.org/0009-0002-2338-5591}{ORCID}\par}
\end{minipage}

\vspace{10pt}
{\color{deepblue}\rule{\textwidth}{1.15pt}}
\vspace{5pt}



% Each paired row starts two sections at precisely the same vertical position.
\pairedsections{
  \sectleft{RESEARCH AREAS}
  {\sffamily\fontsize{9.35}{13.0}\selectfont
  Orthogonal polynomials\par\vspace{3pt}
  Symmetric functions\par\vspace{3pt}
  Integrable systems\par\vspace{3pt}
  Stochastic processes\par}
}{
  \sect{PUBLICATIONS}
  \paper{Approximation of Askey--Wilson roots}
    {J. F. van Diejen, A. Soledispa Tibán, \textbf{A. Vidal}}
    {\emph{Constructive Approximation} (2026). \href{https://doi.org/10.1007/s00365-026-09789-w}{DOI}}
  \paper{The Spectrum of the Finite Open XX Quantum Spin Chain with Transverse Magnetic Boundary Fields via Orthogonal Polynomials}
    {J. F. van Diejen, A. Soledispa Tibán, \textbf{A. Vidal}}
    {\emph{Annales Henri Poincaré} \textbf{27}(4) (2026), 1451--1495. \href{https://doi.org/10.1007/s00023-025-01564-0}{DOI}}
}
\pairedsections{
  \sectleft{EDUCATION}
  \asideentry{PhD in Mathematics}{Universidad de Talca, 2024--present}{Candidate since March 2026. Expected completion: March 2028. Advisor: Jan Felipe van Diejen.}
  \asideentry{M.Sc. in Mathematics}{Universidad de Talca, 2024}{Advisor: Jan Felipe van Diejen.}
  \asideentry{B.Sc. in Mathematics}{Universidad de La Habana, 2020}{}
}{
  \sect{SELECTED PRESENTATIONS}
  % Three most relevant talks from the past two years;
  % exclude meetings held exclusively for students.
  \talk{AUGUST 2026 · COSTA RICA}
    {Hall--Littlewood Polynomials and Interacting Random Walks on Partitions}
    {CIMPA School on Combinatorial Structures in Geometry, Graphs, and Symmetry.}
  \talk{JUNE 2026 · SPAIN}
    {Biased random walks on partitions and hyperoctahedral Hall--Littlewood polynomials}
    {OPSF-S11 Summer School, Universidad de Alcalá.}
  \talk{APRIL 2026 · CHILE}
    {Aproximación de las raíces de los polinomios de Askey--Wilson mediante una iteración de punto fijo}
    {XXXVIII Jornada de Matemática de la Zona Sur, Universidad del Bío-Bío.}
}
\pairedsections{
  \sectleft{ACADEMIC SERVICE}
  {\sffamily\fontsize{9.15}{12.4}\selectfont
  Organizing Committee, \emph{Summer School for Undergraduates in Mathematics} (2026), Universidad de Talca.}
}{
  \sect{TEACHING}
  \lineentry{Lecturer · Universidad de Talca (2025)}
    {Mathematics for Engineering in Nutrition and Dietetics.}
  \lineentry{Teaching Assistant · Universidad de Talca (2022--2024)}
    {Calculus, Linear Algebra, Probability and Statistics, and ODEs.}
  \lineentry{Earlier teaching · Havana (2018--2022)}
    {Universidad de La Habana and CUJAE.}
}

\end{document}